\documentclass[11pt]{article}

%------------------------------------------------
% Packages
%------------------------------------------------
\usepackage[a4paper,margin=2.5cm]{geometry}
\usepackage{amsmath,amssymb,amsthm}
\usepackage{mathtools}
\usepackage{enumitem}
\usepackage{hyperref}

\hypersetup{
    colorlinks=true,
    linkcolor=blue,
    citecolor=blue,
    urlcolor=blue
}

%------------------------------------------------
% Title Information
%------------------------------------------------

\title{\bfseries Adjoint-based Methods for Optimal Control of Hyperbolic
Systems of Conservation Laws}

\author{
Stefan Ulbrich\thanks{Department of Mathematics, Technische Universit\"at
Darmstadt,
\texttt{stefan.ulbrich@tu-darmstadt.de}}
\and
Jannik Breitkopf\thanks{Department of Mathematics, Technische Universit\"at
Darmstadt,
\texttt{jannik.breitkopf@tu-darmstadt.de}}
}

\date{}

%------------------------------------------------

\begin{document}

\maketitle

\vspace{-1em}

\noindent
\textbf{Keywords:}
optimal control, adjoint equation, hyperbolic systems of conservation laws

\vspace{0.5em}


\vspace{1em}

%%%%%%%%%%%%%%%%%%%%%%%%%%%%%%%%%%%%%%%%%%%%%%%%%%
\section*{Abstract}

Hyperbolic systems of conservation problems arise in many applications, for
example in fluid mechanics or traffic flow models. A particular feature lies
in the fact that even for smooth data discontinuities may develop (shocks,
contact discontinuities). This complicates the development of derivative
based optimality conditions and solution methods in the context of optimal
control and inverse problems governed by these PDEs.

We consider strictly hyperbolic $n\times n$ systems of conservation laws in
one space dimension and present a sensitivity calculus for the optimal
control of piecewise $C^1$ solutions containing shocks.
This implies the differentiability of many objective
functionals with respect to controls. To obtain an adjoint-based derivative
representation of the objective functional, we study the adjoint
equation, which is a system of transport equations with discontinuous
coefficients. It turns out that the solution of the adjoint equation is
non-unique and an appropriate boundary condition along shocks is required to
select the correct adjoint state. We present existence, uniqueness and
regularity results for the adjoint equation.

Finally, we discuss challenges and possible approaches to solve the adjoint
equation numerically. Moreover, we derive from the properties of the adjoint
state conclusions for the design of structure-exploiting derivative based methods
for optimal control and inverse problems involving systems of conservation laws.

%%%%%%%%%%%%%%%%%%%%%%%%%%%%%%%%%%%%%%%%%%%%%%%%%%
\vspace{1em}

\noindent
\textbf{Acknowledgements.}
This work was supported by the DFG in the ”Collaborative Research Centre CRC/Transregio 154,
Mathematical Modelling, Simulation and Optimization Using the Example of Gas
Networks”, Project-ID 239904186, Projects A02.


%%%%%%%%%%%%%%%%%%%%%%%%%%%%%%%%%%%%%%%%%%%%%%%%%%
\vspace{1em}

\begin{thebibliography}{9}
\bibitem{BreitkopfUlbrich2025}
J.~Breitkopf and S.~Ulbrich.
\newblock A variational calculus for optimal control of the generalized riemann
  problem for hyperbolic systems of conservation laws, 2025.
\newblock Preprint.
\newblock \href {https://arxiv.org/abs/2509.22076}
  {\nolinkurl{arXiv:2509.22076}}.

\bibitem{BreitkopfUlbrich2026}
J.~Breitkopf and S.~Ulbrich.
\newblock An adjoint calculus for optimal control of entropy solutions of the
  generalized riemann problem for hyperbolic systems of conservation laws with
  shocks, 2026.
\newblock Preprint.
\newblock URL: \url{https://opus4.kobv.de/opus4-trr154/643}.
\end{thebibliography}

\end{document}
